\definecolor{ImperialBlue}{RGB}{0, 0, 255}    % Imperial Blue
\definecolor{ProcessBlue}{RGB}{0, 145, 212}   % Process Blue for contrast
\definecolor{LightBlue}{RGB}{235, 245, 255}   % Adjusted background blue
\definecolor{CoolGrey}{RGB}{158, 161, 162}    % Neutral Grey

\section{High-Level System Design}
The architecture of MERIT is designed to address the ``Black-Box'' nature of traditional
recruitment systems. It is partitioned into three core pipeline stages:
\textbf{Data Extraction and Job Configuration}, \textbf{Multi-Source Evidence Fusion}, and
\textbf{Explainable Interpretation}. Figure \ref{fig:high_level_arch} illustrates how
the recruiter interacts with the three stages of MERIT.

% Figure 4.1: High-Level System Design
\begin{figure}[htbp]
    \centering
    \begin{tikzpicture}[
        node distance=1.2cm and 1.8cm,
        arch_domain/.style={draw=ImperialBlue, fill=LightBlue!30, dashed, rounded corners, thick, inner sep=15pt},
        arch_logic/.style={rectangle, draw=ImperialBlue, fill=white, text width=2.8cm, align=center, minimum height=1.1cm, rounded corners, thick, font=\footnotesize\sffamily, drop shadow},
        arch_user/.style={ellipse, draw=ImperialBlue, fill=ImperialBlue!10, minimum width=1.8cm, minimum height=0.8cm, font=\small\bfseries\sffamily, thick},
        arrow/.style={-{Stealth[scale=1.0]}, thick, draw=ImperialBlue},
        domain_label/.style={font=\footnotesize\bfseries\sffamily, ImperialBlue, anchor=center}
    ]
        % I. Extraction
        \node [arch_logic] (db1) {Supabase\\Unified Store};
        \node [arch_logic, above=0.8cm of db1] (ext) {Multi-source\\Parsing Engine};
        
        % II. Evidence Fusion
        \node [arch_logic, right=2.2cm of db1] (fusion) {Bayesian Fusion\\\& Scoring Engine};
        
        % III. Explainable Interpretation
        \node [arch_logic, right=2.2cm of fusion, yshift=1.6cm] (run) {Candidate Ranker\\\& Scoring Pipeline};
        \node [arch_logic, below=0.5cm of run] (audit) {ShaRP Explanation\\\& Audit Trails};

        % Containers (Using FIT to ensure boxes are inside)
        \begin{scope}[on background layer]
            \node [arch_domain, fit=(ext) (db1)] (area1) {};
            \node [arch_domain, fit=(fusion)] (area2) {};
            \node [arch_domain, fit=(run) (audit)] (area3) {};
        \end{scope}

        % User (Recruiter) - Increased height for absolute clearance
        \node [arch_user, above=5.5cm of area2] (u2) {Recruiter};
        
        % Flows (Segmented for precise label control)
        % To Extraction
        \draw [arrow] (u2.west) -- (u2.west -| area1.center) node[midway, above, font=\footnotesize\sffamily, black, fill=white, inner sep=1pt] {Ingest Documents} -- (ext.north);
        
        % To Configuration
        \draw [arrow] (u2.south) -- (fusion.north) node[pos=0.3, left, font=\footnotesize\sffamily, black] {Set Weights};
        
        % Internal Flows
        \draw [arrow] (ext) -- (db1);
        \draw [arrow] (db1.east) -- (fusion.west);
        \draw [arrow] (fusion.east) -- (run.west);
        \draw [arrow] (run) -- (audit);
        
        % Feedback Loop (Mirrored Symmetry)
        \draw [arrow] (audit.east) -- ++(1.5,0) |- (u2.east) node[pos=0.75, above, font=\footnotesize\sffamily, black, fill=white, inner sep=1pt] {Explainable Insights};

        % Domain Headers (Drawn last with white background to create "gap" in arrows)
        \node [domain_label, fill=white, inner sep=2pt] at (area1.north) {I. DATA EXTRACTION};
        \node [domain_label, fill=white, inner sep=2pt] at (area2.north) {II. MULTI-SOURCE FUSION};
        \node [domain_label, fill=white, inner sep=2pt] at (area3.north) {III. RANKING \& EXPLAINABILITY};

    \end{tikzpicture}
    \caption[MERIT three-stage architecture]{High-Level System Design: The three stages of MERIT: Ingestion, Scoring, Explainability}
    \label{fig:high_level_arch}
\end{figure}


\subsection{Rationale}
MERIT employs a decoupled client-server architecture to separate the
concerns of data ingestion, job configuration, and the presentation of results.
We used Flask (Python) as the backend for data ingestion and fusion due to its rich ecosystem
of Python libraries such as NLTK, spaCy, and Scikit-learn.
To ensure fast prototyping and development, we used TypeScript for the frontend, 
with React and Next.js for the UI. Note that our choice of tech stack is not mandatory, it is 
just a recommendation for fast, scalable and maintainable development.

A relational PostgreSQL store (via Supabase) was selected to enforce data integrity,
ensuring that every ranked result is linked via foreign keys to a specific,
immutable recruiter configuration.

\subsection{Reliability \& Scalability}
While prototyping MERIT, we encountered many major obstacles
during the development of the system. Below are the two 
greatest non-pillar challenges that we faced:

\subsubsection{The N+1 Query Bottleneck}
A significant scalability challenge was the \textbf{N+1 query bottleneck}. When
retrieving $N$ candidates, the backend was previously making $N$ subsequent database
calls to fetch associated skills and scores. By leveraging PostgreSQL's relational
embedding capabilities via the Supabase client, we transitioned query complexity from
$O(N)$ to $O(1)$. The benefits of this runtime optimisation are significant when considering
a large candidate pool, which will be almost always the case in real-world
applications.

\subsubsection{Data Integrity \& Pipeline Robustness}
One of the most significant challenges we faced was in the parsing of unstructured data (CVs and JDs),
where the inherent unpredictability of the extracted text could cause pipeline failures.
Our approach to overcome this was the introduction of strict JSON schema validation at the ingestion boundary, ensuring
that only verified signals enter the multi-source fusion engine. We followed this up with global Flask middleware to
intercept and handle any exceptions throughout the backend pipeline. This helped ensure the overall integrity and reliability
of the explainability layer.
